\section{ESP32 GPIO Pin Configuration and Function Mapping}

Complete reference for all GPIO pins used in NexusFlow hardware variants.

\subsection{Consolidated Pin Mapping Table}

\begin{table}[H]
\centering
\caption{Complete ESP32 GPIO Pinout Reference}
\label{tbl:gpio_reference}
\begin{tabular}{|c|c|c|c|c|}
\hline
\textbf{GPIO} & \textbf{Function} & \textbf{I/O Type} & \textbf{4-Ch} & \textbf{6-Ch/8-Ch} \\
\hline
0 & Mode/Pairing Button & INPUT & ✓ & ✓ \\
\hline
4 & Red Status LED & OUTPUT & ✓ & ✓ \\
\hline
5 & Relay 3 Control & OUTPUT & ✓ & ✓ \\
\hline
12 & Touch Sensor 6 Input & INPUT & ✗ & ✓ \\
\hline
13 & Touch Sensor 5 Input & INPUT & ✗ & ✓ \\
\hline
14 & Touch Sensor 4 Input & INPUT & ✗ & ✓ \\
\hline
15 & Touch Sensor 3 Input & INPUT & ✗ & ✓ \\
\hline
16 & Relay 5 Control & OUTPUT & ✗ & ✓ \\
\hline
17 & Relay 4 Control & OUTPUT & ✗ & ✓ \\
\hline
18 & Relay 2 Control & OUTPUT & ✓ & ✓ \\
\hline
19 & Relay 1 Control & OUTPUT & ✓ & ✓ \\
\hline
21 & I2C SDA (Sensors) & I2C & ✓ & ✓ \\
\hline
22 & I2C SCL (Sensors) & I2C & ✓ & ✓ \\
\hline
23 & Green Status LED & OUTPUT & ✓ & ✓ \\
\hline
25 & Relay 6 Control & OUTPUT & ✗ & ✓ \\
\hline
27 & Mode-Select Slider & ADC INPUT & ✓ & ✓ \\
\hline
32 & Touch Sensor 1 Input & INPUT & ✓ & ✓ \\
\hline
33 & Touch Sensor 2 Input & INPUT & ✓ & ✓ \\
\hline
\end{tabular}
\end{table}

\subsection{4-Channel Variant Configuration}

For 4-relay systems (e.g., basic home setup):

\begin{lstlisting}[language=Kotlin, caption=4-Channel GPIO Configuration, label=lst:gpio_4ch]
object HardwareConfig4Ch {
    const val RELAY_COUNT = 4
    const val TOUCH_COUNT = 4
    
    val RELAY_PINS = intArrayOf(
        19, // Relay 1
        18, // Relay 2
        5,  // Relay 3
        17  // Relay 4
    )
    
    val TOUCH_PINS = intArrayOf(
        32, // Touch 1
        33, // Touch 2
        14, // Touch 3
        15  // Touch 4
    )
    
    const val I2C_SDA = 21
    const val I2C_SCL = 22
    const val GREEN_LED = 23
    const val RED_LED = 4
    const val MODE_BUTTON = 0
    const val MODE_SLIDER = 27
}
\end{lstlisting}

\subsection{6-Channel Variant Configuration}

For 6-relay systems (typical household):

\begin{lstlisting}[language=Kotlin, caption=6-Channel GPIO Configuration, label=lst:gpio_6ch]
object HardwareConfig6Ch {
    const val RELAY_COUNT = 6
    const val TOUCH_COUNT = 6
    
    val RELAY_PINS = intArrayOf(
        19, // Relay 1
        18, // Relay 2
        5,  // Relay 3
        17, // Relay 4
        16, // Relay 5
        25  // Relay 6
    )
    
    val TOUCH_PINS = intArrayOf(
        32, // Touch 1
        33, // Touch 2
        15, // Touch 3
        14, // Touch 4
        13, // Touch 5
        12  // Touch 6
    )
    
    const val I2C_SDA = 21
    const val I2C_SCL = 22
    const val GREEN_LED = 23
    const val RED_LED = 4
    const val MODE_BUTTON = 0
    const val MODE_SLIDER = 27
}
\end{lstlisting}

\subsection{8-Channel Variant Configuration}

For 8-relay systems (advanced installations):

\begin{lstlisting}[language=Kotlin, caption=8-Channel GPIO Configuration, label=lst:gpio_8ch]
object HardwareConfig8Ch {
    const val RELAY_COUNT = 8
    const val TOUCH_COUNT = 6  // Limited by available pins
    
    val RELAY_PINS = intArrayOf(
        19, // Relay 1
        18, // Relay 2
        5,  // Relay 3
        17, // Relay 4
        16, // Relay 5
        25, // Relay 6
        27, // Relay 7 (shared with mode slider, configurable)
        2   // Relay 8 (additional GPIO if available)
    )
    
    val TOUCH_PINS = intArrayOf(
        32, // Touch 1
        33, // Touch 2
        15, // Touch 3
        14, // Touch 4
        13, // Touch 5
        12  // Touch 6
    )
    
    const val I2C_SDA = 21
    const val I2C_SCL = 22
    const val GREEN_LED = 23
    const val RED_LED = 4
    const val MODE_BUTTON = 0
}
\end{lstlisting}

\section{Critical GPIO Warnings}

\subsection{Strapping Pins}

\begin{table}[H]
\centering
\caption{ESP32 Strapping Pins and Boot Behavior}
\label{tbl:strapping_pins}
\begin{tabular}{|c|c|c|l|}
\hline
\textbf{GPIO} & \textbf{Strapping Function} & \textbf{Default} & \textbf{Issue} \\
\hline
0 & JTAG/Download mode & Pull-up & LOW triggers flash mode \\
\hline
2 & Internal VREF (analog) & Pull-down & Not recommended for GPIO \\
\hline
5 & SPI FLASH Mode & Pull-up & Interferes with SPI boot \\
\hline
12 & SPI FLASH Voltage & Pull-down & HIGH triggers incorrect voltage \\
\hline
15 & SDIO Slave (SD card) & Pull-up & LOW enables SD mode \\
\hline
\end{tabular}
\end{table}

\textbf{Implications for NexusFlow:}

\begin{itemize}
    \item \textbf{GPIO 0:} Used as mode/pairing button. Must include debouncing firmware to prevent accidental boot mode trigger during normal operation.
    
    \item \textbf{GPIO 12, 15:} Used for touch inputs in 6/8-channel variants. Small boot sensitivity; firmware debouncing required. If unexpected resets occur, consider PCB redesign with dedicated GPIO or resistive debouncing.
    
    \item \textbf{GPIO 5:} Used for Relay 3. SPI FLASH mode strapping; ensure proper pull-up configuration. Generally safe after boot, but avoid holding LOW during power-up.
\end{itemize}

\subsection{Power Supply Considerations}

\begin{itemize}
    \item \textbf{3.3V Rail:} Supplies ESP32, TTP223 modules, AHT10 sensor, status LEDs
    \item \textbf{5V Rail:} Supplies relay module coils only
    \item \textbf{Ground:} Common ground for all components; use star grounding for noise immunity
\end{itemize}

\textbf{Current Budget:}
\begin{itemize}
    \item ESP32 Core: ~80 mA typical
    \item 6 Relays (coil, 5V): ~120 mA (20 mA each)
    \item TTP223 × 6: ~6 mA
    \item AHT10: ~0.3 mA
    \item Status LEDs: ~20 mA
    \item \textbf{Total 3.3V:} ~150 mA
    \item \textbf{Total 5V:} ~120 mA
\end{itemize}

A 1A @ 5V USB power supply is sufficient for development. Production systems should use a dedicated power supply rated for 2A @ 5V to accommodate transient relay switching current spikes.

\section{GPIO Electrical Characteristics}

\begin{table}[H]
\centering
\caption{ESP32 GPIO Electrical Specifications}
\label{tbl:gpio_specs}
\begin{tabular}{|l|c|c|}
\hline
\textbf{Parameter} & \textbf{Min} & \textbf{Max} \\
\hline
Output High Voltage (IOH = -1 mA) & 2.4V & 3.3V \\
\hline
Output Low Voltage (IOL = 10 mA) & 0V & 0.4V \\
\hline
Input High Voltage & 2.0V & 3.3V \\
\hline
Input Low Voltage & 0V & 0.8V \\
\hline
Output Drive Capability & 12 mA & \text{(per GPIO)} \\
\hline
Input Capacitance & 0 pF & 150 pF \\
\hline
\end{tabular}
\end{table}

\textbf{Design Implications:}

\begin{itemize}
    \item GPIO outputs 3.3V—safe for TTP223 input (rated 3.6V max)
    \item GPIO outputs drive relay module transistors directly (10–20 mA requirement, GPIO capable of 12 mA)
    \item 5V relay module outputs must NOT connect to GPIO (overvoltage risk)
    \item Pull-up/pull-down resistors: 4.7 kΩ recommended for I2C (per I2C spec)
\end{itemize}

\section{Firmware GPIO Initialization}

\begin{lstlisting}[language=Kotlin, caption=GPIO Initialization Code, label=lst:gpio_init]
void initializeGPIO() {
    // Configure relay pins as outputs
    for (int relay : RELAY_PINS) {
        pinMode(relay, OUTPUT);
        digitalWrite(relay, HIGH);  // Relays active-LOW; start inactive
    }
    
    // Configure touch sensor pins as inputs
    for (int touch : TOUCH_PINS) {
        pinMode(touch, INPUT);
    }
    
    // Configure I2C pins (handled by I2C library)
    // Wire.begin(I2C_SDA, I2C_SCL);
    
    // Configure status LEDs
    pinMode(GREEN_LED, OUTPUT);
    pinMode(RED_LED, OUTPUT);
    digitalWrite(GREEN_LED, LOW);   // Start off
    digitalWrite(RED_LED, LOW);     // Start off
    
    // Configure mode button
    pinMode(MODE_BUTTON, INPUT_PULLUP);
    
    // Configure mode slider (ADC)
    pinMode(MODE_SLIDER, INPUT);
    
    Serial.println("GPIO initialization complete");
}

// Helper functions for relay control
void turnRelayOn(int relayIndex) {
    if (relayIndex >= 0 && relayIndex < RELAY_COUNT) {
        digitalWrite(RELAY_PINS[relayIndex], LOW);  // Active-LOW
    }
}

void turnRelayOff(int relayIndex) {
    if (relayIndex >= 0 && relayIndex < RELAY_COUNT) {
        digitalWrite(RELAY_PINS[relayIndex], HIGH);  // Inactive
    }
}

bool isRelayOn(int relayIndex) {
    if (relayIndex >= 0 && relayIndex < RELAY_COUNT) {
        return digitalRead(RELAY_PINS[relayIndex]) == LOW;
    }
    return false;
}
\end{lstlisting}

\section{Troubleshooting GPIO Issues}

\subsection{Relay Not Switching}

\begin{enumerate}
    \item Verify GPIO pin correct in configuration (matches your hardware variant)
    \item Check relay module power (5V present on VCC)
    \item Verify active-LOW logic: \texttt{digitalWrite(pin, LOW)} should energize relay
    \item Check current limiting: relay module should have built-in transistor driver
    \item Measure GPIO voltage: should toggle between 0V (ON) and 3.3V (OFF)
\end{enumerate}

\subsection{Touch Input Not Responding}

\begin{enumerate}
    \item Verify TTP223 powered from 3.3V (NOT 5V)
    \item Confirm output pin connected correctly to GPIO
    \item Check GPIO voltage: should oscillate between 0V and 3.3V when touched
    \item Test with Serial output: \texttt{Serial.println(digitalRead(touchPin))}
    \item Verify debouncing delay (TTP223 spec 60 ms)
\end{enumerate}

\subsection{I2C Sensor Communication Failure}

\begin{enumerate}
    \item Verify SDA on GPIO 21, SCL on GPIO 22
    \item Check 4.7 kΩ pull-up resistors present on both lines
    \item Confirm AHT10 at I2C address 0x38
    \item Oscilloscope test: verify rising/falling edges clean (not rounded)
    \item Try reducing I2C clock speed if using longer wires (e.g., 100 kHz instead of 400 kHz)
\end{enumerate}

\subsection{Unexpected Resets or Boot Failures}

\begin{enumerate}
    \item Check strapping pin state (GPIO 0, 12, 15) during power-up
    \item Add debouncing on GPIO 0 (mode button) to prevent accidental flash mode
    \item Verify adequate power supply (1A @ 5V minimum for 6 relays)
    \item Check USB cable quality—flaky cables cause brownouts
    \item Monitor serial output during boot: watch for error messages
    \item If GPIO 12 or 15 issues persist, consider redesign to use alternative pins
\end{enumerate}

\section{GPIO Pin Assignment Rationale}

\subsection{Why These Pins?}

\begin{itemize}
    \item \textbf{Relays (19, 18, 5, 17, 16, 25):} Contiguous logic outputs; easy to manage in arrays
    \item \textbf{Touch (32, 33, 14, 13, 12, 15):} Separate from relay pins; avoids cross-talk
    \item \textbf{I2C (21, 22):} Dedicated I2C pins on ESP32; avoid potential conflicts
    \item \textbf{Status LEDs (23, 4):} Unused pins; easily accessible for debugging
    \item \textbf{Mode Button (0):} Special function but acceptable with firmware debouncing
\end{itemize}

\subsection{Alternative Pin Assignments}

If conflicts arise with other modules:

\begin{itemize}
    \item \textbf{Additional Relays:} GPIO 2, 26, 27 available (but check prior use)
    \item \textbf{Additional Sensors:} SPI pins 23, 18, 19, 5 available if SPI not used
    \item \textbf{ADC inputs:} GPIO 34, 35, 36, 39 dedicated analog inputs (no output capability)
\end{itemize}

This appendix serves as definitive GPIO reference during hardware assembly, firmware development, and troubleshooting phases.
